% ═══════════════════════════════════════════════════════════
% §3 Discussion
% ═══════════════════════════════════════════════════════════
\section{Discussion}
\label{sec:discussion}

\subsection{Summary of findings}
This work establishes three results that advance our understanding of how
discrete mathematical structures interact with protein coevolution analysis.
First, we demonstrate that Quine--McCluskey-minimised Karnaugh-map circuits,
learned from consensus-residue identities in training proteins, transfer to
held-out proteins at $5.1\times$ random enrichment for native contacts.
Second, we show that the cell-rate prior derived from these circuits improves
continuous DCA ranking when fused via rank aggregation, providing the first
evidence that Boolean-minimisation features carry structural information
orthogonal to standard covariation analysis.  Third, we confirm at proper
statistical power ($p < 10^{-77}$) that Gray-code adjacency is enriched among
native contacts and demonstrate through permutation testing that this
enrichment is specific to the He/Gray encoding rather than an artefact of
generic chemistry clustering.

\subsection{Comparison with existing methods}
Our circuits achieve prec@L/5 = 0.179 on held-out proteins, which is
substantially below PSICOV's sparse inverse covariance method (0.727) but
above plain mutual information without correction (0.097).  The gap reflects
the fundamental difference between discrete consensus-identity encoding with
Boolean minimisation and continuous covariance-based inference with
phylogenetic correction.  These approaches are complementary rather than
competing: the rank-fusion experiment demonstrates that combining them yields
better predictions than either alone.

The biological interpretation of learned circuit rules is a key advantage
over neural-network-based methods.  Our Level-1 essential implicants encode
hydrophobic/aromatic core packing and sulfur-mediated contacts in directly
readable sum-of-products form.  At Level-2 resolution, specific residue-pair
vocabularies such as $\{M,F\}\times\{M,F,Y,W\}$ correspond to known
aliphatic/aromatic packing interactions, rediscovered from sequence data
without biophysical prior knowledge.

\subsection{Limitations}
Several limitations constrain the scope of our conclusions.  Consensus-residue
features collapse within-column frequency variation into a single code per
position, discarding information that richer encodings might exploit.
Separation bins use globally defined boundaries without per-protein length
normalisation.  The cross-dataset generalisation experiment on EVmutation was
inconclusive because correct parsing of the A2M alignment format requires
handling insert-state columns that do not map linearly to PDB residue numbers;
this is a well-known challenge of the EVmutation format rather than a flaw in
our framework.  Finally, the iterative direct-information computation from
our mean-field implementation remains unreliable on families with many fully
conserved columns due to pseudo-inverse conditioning; we adopted APC-corrected
Frobenius norm as the primary DCA score for this reason.

\subsection{Future directions}
Three extensions follow naturally.  First, replacing binary conservation flags
with continuous entropy bins as additional Karnaugh-map dimensions could
capture finer-grained structural information about the buried core.  Second,
integrating circuit-derived rules as input features to deep learning models
would combine the interpretability of Boolean logic with the representational
power of attention-based architectures.  Third, applying the complete pipeline
to CASP free-modelling targets would test whether the observed transfer
generalises to harder targets with shallower alignments and more diverse
structural classes.
